\subsection*{哪些检查发生在受理以前}

恢复协议必须先能构造一个可信的 S 和一份可解释输入。入口先认证调用者，核对其账户权限，验证动作与操作许可，检查会话状态与时间，解析请求体并检查字段、类型和长度。JSON 畸形或重复成员在本版中返回 \code{400 invalid\_input}；可以解析但字段不合规则返回 \code{422 invalid\_input}，例如未知字段或 quantity 为字符串。两类错误都不插入购买身份。认证失败也不向 \code{purchases} 写一条该账户的拒绝记录；服务尚未确认来者有资格占用这个空间。

这条边界防止攻击者拿别人的 K1 把身份提前填成认证失败，也使一次截断报文不会永久占住正确购买。请求格式修正以后，调用者可按规则继续原 K1。不过 \code{invalid\_input} 只说明本次尝试没有受理，不保证别的同 S 请求从未到达。若另一个完整 R1 已提交，修正后的 R2 仍须查找原记录并核对输入，不能把本次校验失败当作全局无记录证明。

通过前置检查以后，服务进入数据库协议，尝试插入完整 S、输入版本与规范化输入。只有实际插入者进行首次业务判断。取得身份之后，余额不足或经明确规则判断的报价过期，可以保存为终态；数据库不可用、锁等待超时、无法取得报价依据等运行故障则不自动成为业务拒绝。本章报价来自当前事务中可验证的本地依据，故没有把远程报价超时藏进终态判断。

\begin{figure}[htbp]
\centering
\begin{tikzpicture}[x=1cm,y=1cm]
\node[booknode] (a) at (0,0) {认证、格式与受理检查};
\node[booknote,text width=3.3cm] (e) at (6,0) {前置失败：本次未受理\\不能据此判断其他同 S 尝试};
\node[booknode] (b) at (0,-1.4) {完整 S 竞争并核对原输入};
\node[booknote,text width=3.3cm] (f) at (6,-1.4) {同 S 异参：409\\保留已有绑定与终态};
\node[bookresult] (c) at (0,-2.8) {首次取得身份，执行事务};
\node[booknode] (g) at (-2,-4.2) {成功：订单与扣款};
\node[booknode] (h) at (3,-4.2) {终态拒绝：保存原因};
\draw[bookarrow] (a)--(b);
\draw[bookarrow] (a)--(e);
\draw[bookarrow] (b)--(f);
\draw[bookarrow] (b)--(c);
\draw[bookarrow] (c)--(g);
\draw[bookarrow] (c)--(h);
\end{tikzpicture}
\caption{前置拒绝、同身份异参冲突与业务终态拒绝分别发生在不同步骤。}
\label{fig:ch05-admission-branches}
\end{figure}


同一字段还可能有两类错误。\code{quote\_id} 的类型必须是字符串，这属于输入格式；Q7 是否属于本账户、是否允许该商品成交，涉及权限与业务规则。服务应先检查访问安全，避免通过恢复接口泄露不该查看的报价；随后按已说明的首次成交规则判断。对于已有同参终态，Q7 后来自然过期不再否定旧成交。需要当前权限才能访问的限制仍继续执行，但不借这项限制改写 P1 的资金事实。

\subsection*{余额不足怎样成为最后结论}

另取初始状态：账户余额 2000，购买、订单与流水为空，S1 请求 Q7 的 3000 分。A 插入 S1，条件 UPDATE 检查 \code{balance\_minor>=3000}，返回零行。账户存在且权限已核实，所以本例明确得出余额不足。A 保存 \code{P1.status=insufficient\_funds}，记录原输入、Q7、CNY、所需金额 3000 与拒绝原因。在这段历史中，拒绝提交以前没有其他账户变化；提交后没有订单和扣款，余额仍为 2000。

购买记录把这个终态命名为 \code{insufficient\_funds}。响应将 status 设为 \code{rejected}，将 reason 设为 \code{insufficient\_funds}。内部终态与外部表示的映射由 \code{purchase-result-v1} 固定。响应正文必须保留 P1 与原绑定输入所需字段，不能只返回一句供界面展示的余额不足提示语。reason 是机器可判断的原因码，人类提示语可按当前语言生成；若提示语包含具体成交数值，应仍来自保存事实。

本例不在拒绝结果中返回 \code{available\_balance\_minor}。条件 UPDATE 零行本身只提供条件不满足的结论，服务后来另查余额，可能已经看到一笔并发充值。把那个值当成判断时余额，会混入新的时点。若产品确实需要拒绝时精确余额，可先取得账户行锁并在受保护读取中判断，或者采用另一套能记录判定值的受控语句；这会改变并发路径，需要单独核对。本章拒绝结果只声称当时所需 3000 的扣款条件未成立。

\begin{BookCode}
{"result_version":"purchase-result-v1",
 "purchase_id":"P1","status":"rejected",
 "reason":"insufficient_funds",
 "product_id":"paper-001","quantity":1,
 "quote_id":"Q7","currency":"CNY",
 "required_amount_minor":3000}
\end{BookCode}

后来 Alice 充值 4000，账户为 6000。R2 恢复 S1，仍返回这份拒绝，不产生 O1、D1。Alice 此时确认重新购买，应使用新的 S2、保存新输入并按届时有效报价提交。若原 Q7 仍可用，新操作可以采用它；若已经过期，则接受新报价。新键表达新的购买确认，旧键则继续解释充值以前已经结束的操作。

\begin{figure}[htbp]
\centering
\begin{tikzpicture}[x=1cm,y=1cm]
\node[booknode,text width=3cm] (a) at (0,0) {余额 2000\\S1 需扣 3000};
\node[bookresult,text width=3cm] (b) at (4,0) {S1 拒绝提交\\无订单与扣款};
\node[booknode,text width=3cm] (c) at (8,0) {之后充值 4000\\余额变为 6000};
\draw[bookarrow] (a)--(b);
\draw[bookarrow] (b)--(c);
\node[booknode,text width=4cm] (d) at (2,-2) {恢复原 S1：仍返回原拒绝};
\node[booknode,text width=4cm] (e) at (8,-2) {新购买 S2：重新判断\\扣 3000 后余额 3000};
\draw[bookalert] (b)--(d);
\draw[bookarrow] (c)--(e);
\end{tikzpicture}
\caption{本例充值发生在拒绝提交之后；新增购买使用新身份重新判断业务条件。}
\label{fig:ch05-rejection-recharge}
\end{figure}


这段历史把充值安排在原拒绝提交之后，检查的是旧结论怎样恢复。充值与余额不足判断并发时，判定和提交可以观察到不同余额，第 6.2 节将单独检查其生效位置。

首次余额不足继续沿用前四章的 409 Conflict。响应体中的 \code{status} 取 \code{rejected}，\code{reason} 取 \code{insufficient\_funds}，说明拒绝终态已经提交。恢复到这份终态时返回 200，说明成功取回原结论。409 也用于同身份异参冲突，客户端因此还要读取业务码与本次受理信息。422 仅用于本版的字段语义校验错误；HTTP 状态类别提供通信层概括，是否保存终态由响应体和接口规则说明\cite{http}。

\subsection*{临时故障保留原身份恢复}

暂时故障可以发生在取得 S 以前、事务内部、发送 COMMIT 时或结果返回途中。如果数据库明确报告整笔事务回滚，本次尝试没有提交其修改；其他同 S 请求仍可能接着提交。如果提交确认丢失，则连本次事务是否成立也未知。两种情况下，客户端都保留原 S，恢复者依旧通过唯一仲裁与终态读取继续，区别在于服务内部已经掌握多少证据。

本接口可对暂时无法完成或确认的尝试返回 503 Service Unavailable，附带 Retry-After 建议何时再试。Retry-After 是 HTTP 中表达建议等待时间或时点的响应头；它用于减轻即时重试压力，不担保到期时一定有结果。响应体可以写 \code{code=result\_unconfirmed}、\code{operation\_outcome=unknown}、\code{retry\_same\_identity=true}，要求客户端保留 E1、K1 和原输入。这里 unknown 描述本次观察，数据库中可能已经有成功终态\cite{http}。

服务确知本次事务回滚时，可以在诊断信息中记录 \code{attempt\_rolled\_back}；不要由此给整项 S 返回一个可永久重放的 \code{failed}。另一同身份事务可能成功，或者后续原身份重试才第一次成功。除非业务明确决定终止该操作并以共同规则持久保存终态，否则临时错误响应不写入原结果。本例服务不会把磁盘暂时不可用或内部程序异常保存成 K1 的永久业务拒绝。

数据库中的记录若处于意外的已提交 \code{pending}，也属于设计前提被破坏。普通查询看到它时，不能直接扣款来猜测前一个进程做到了哪步。应返回明确的待处理错误并转入受控修复。当前协议从未提交 \code{pending}，所以也没有可靠的外部 \code{processing} 状态可用于 202 Accepted。202 适合已有持久受理证据且有后续完成安排的接口，本章没有把它当作一个万能等待响应。

\subsection*{冲突保留原记录的含义}

设 S 已绑定 \code{paper-001}、数量 1、Q7，终态成功。R2 却携带 \code{paper-002} 或 Q8。服务先按完整 S 取得原记录，按保存的 \code{normalization\_version} 整理本次输入，再比较 \code{canonical\_input}。不同就返回 409 Conflict，业务码 \code{idempotency\_input\_conflict}。响应可以提供允许查看的原 \code{purchase\_id=P1} 和冲突字段名，但不能把 \code{paper-001} 的成功结果包装成 \code{paper-002} 的成功响应。

冲突结果属于 R2 这次请求与原绑定之间的关系，不覆盖 \code{P1.result\_payload}，也不给 S 写第二个终态。同一 S 后来携带正确原输入，仍能恢复 P1。否则一个误填字段的重试就能毁掉合法购买的恢复能力。原终态若是余额不足，处理完全相同：先比较输入，异参则冲突；同参再返回原拒绝。为拒绝也保存完整身份与输入，正是为了让这条路径有明确依据。

尚无旧终态的同 S 异参竞争，沿用第 4.3 节的首次提交绑定规则。调用方仍需保持原输入；服务在输入混用时继续限制一份终态、保持绑定不可覆盖，并向另一份输入报告冲突。

\subsection*{把接口错误与持久事实对齐}

可以按是否取得持久终态安排响应。首次成功为 201，首次明确余额不足为 409；同参恢复，无论原成功或原拒绝，均用 200 携带原结果。前置错误返回 \code{invalid\_input} 或认证授权错误，不写购买身份；异参为 409；当前未查到为 404；暂时无法取得可信结论为 503。客户端管理本地操作时，要同时看业务码、完整身份和已保存结果，不能用所有非 2xx 都换新键的通用逻辑。

\begin{table}[htbp]
\centering\small
\begin{tabularx}{\linewidth}{@{}p{4.2cm}XX@{}}
\toprule
本次响应 & 关于原 S 已知的事实 & 后续动作\\
\midrule
201 succeeded & 成功终态已提交 & 保存 P1 及原结果\\
409 rejected 与业务原因 & 拒绝终态已提交 & 结束该操作，新确认用新键\\
200 原成功或原拒绝 & 读取到已保存终态 & 按原结论更新本地记录\\
409 input\_conflict & S 已绑定不同输入 & 保留旧记录，核对误复用\\
404 operation\_not\_found & 本次查询未找到终态 & 保留原 S，查询或同参重试\\
503 result\_unconfirmed & 当前未能确认结论 & 保留原 S，按建议恢复\\
\bottomrule
\end{tabularx}
\caption{通信响应、原操作事实与客户端后续动作分别对应。}
\end{table}

前置 \code{invalid\_input} 若明确报告 \code{admission=not\_admitted}，这个字段也只描述本次请求没有进入购买事务。服务若提供 accepted、replayed 或 unknown 等受理说明，必须定义其证据来源；不要返回一个笼统 \code{accepted=false} 让客户端误以为同一 S 在全系统从未受理。已完成结果与当前请求的处理状态可以同时出现，例如 R2 因权限失效未被允许读取，但 P1 的成功事实仍存在。

E1 到期后，旧操作入口按第三章的关闭边界返回 \code{410 operation\_expired}。410 说明恢复资格结束，原购买仍可能已经成功，客户端应沿账户内的订单和购买事实核查。截止前已受理的执行继续保有身份与结果直到安全终结；业务档案可以保留更久，其期限与旧 POST 的准入期限分别规定。
\FloatBarrier
